\section{논의}
\label{sec:discussion}

본 논문에서 제시한 두 연구는 UAV 상태 추정 문제의 서로 다른 수준을 다루며,
각 결과는 교차 시사점을 갖는다.

\textbf{임베디드 비행 컨트롤러를 위한 필터 선택.}
3절의 IMU 비교 결과, ESKF가 시뮬레이션에서 가장 낮은 자세 RMSE를 달성하였으나
아두이노급 하드웨어에서 사이클당 연산 비용이 가장 높다.
PX4와 ArduPilot은 \SI{480}{\mega\hertz}로 동작하는 STM32H7급 프로세서에서 EKF를 실행하므로
23상태 EKF의 추가 연산 비용이 허용 가능하다.
더 제한된 자원을 가진 마이크로컨트롤러의 경우 Mahony 필터가 실용적인 대안을 제공한다:
적분 채널을 통한 바이어스 보정 기능이 EKF의 바이어스 상태와 유사하면서도
연산 예산의 일부만을 필요로 한다.
반면 상보 필터는 바이어스 보정이 없어 자이로스코프 드리프트가 누적되는
장시간 비행에는 권장되지 않는다.

\textbf{지터가 필터 정확도에 미치는 영향.}
3절 A항의 지터 분석 결과, 아두이노 Mega의 I\textsuperscript{2}C 샘플링은
% TODO: 측정된 지터 값 삽입
수준의 샘플 간 타이밍 편차를 유발함을 확인하였다.
EKF와 ESKF에서 지터는 $\Delta t$를 통한 과정 노이즈 공분산 스케일링을 직접 교란시키므로,
이를 무시하면 예측 공분산에 체계적 바이어스가 발생한다.
PX4와 ArduPilot의 비행 컨트롤러 구현은 하드웨어 인터럽트 수준에서 IMU 샘플에 타임스탬프를 부여하고
샘플별 정확한 $\Delta t$를 전파함으로써 이 문제를 해결하며,
이는 확률적 필터의 임베디드 구현에서 반드시 적용해야 할 관행이다.

\textbf{PPID 이득 민감도와 추정기 결합.}
4절의 MuJoCo 시뮬레이션 결과, PPID 자세 추적 성능이
내부 루프~\eqref{eq:rate_loop}에 입력되는 각속도 추정값 $\hat{\boldsymbol{\omega}}$의 정확도에
민감하게 반응함을 확인하였다.
3절 A항에서 특성화한 자이로스코프 노이즈 수준이 높을수록 $\hat{\boldsymbol{\omega}}$가 저하되어
내부 루프 각속도 오차가 증가하고, 이는 제어기 대역폭 이상의 주파수에서
자세 응답의 극한 사이클 진동으로 나타난다.
이러한 결합 관계는 PX4와 ArduPilot 모두에 존재하는 각속도 측정값의 노치 필터 적용을 정당화하며,
실용적인 배포에서 PPID 이득 튜닝 이전에 IMU 노이즈 특성화가 선행되어야 함을 시사한다.

\textbf{한계.}
본 연구의 주요 한계는 하드웨어 실험 중 교정된 외부 기준 자세 센서가 없어
물리적 플랫폼에서의 필터 정확도 정량적 검증이 불가능하다는 점이다.
또한 MuJoCo 시뮬레이션이 실제 SK450 기체 모델 대신 Skydio~2 모델을 사용하여
공력 파라미터 불일치로 인해 시뮬레이션-실물 간 격차가 정량화되지 않은 채 존재한다.

\section{결론}
\label{sec:conclusion}

본 논문은 UAV 상태 추정에 관한 두 가지 상호 보완적인 연구를 제시하였다.
첫 번째 연구는 MPU-9250/아두이노 하드웨어 플랫폼에서 여섯 가지 자세 추정 필터를 구현하고 비교하였으며,
센서 노이즈와 지터를 특성화하고 MuJoCo 시뮬레이션 기준값 대비 필터 정확도를 평가하였다.
ESKF가 가장 낮은 RMSE를 달성하였고, Mahony 필터는 자원이 제한된 임베디드 배포에서
가장 우수한 정확도-비용 비율을 제공하였다.
두 번째 연구는 SK450 + CUAV 7 Nano 하드웨어 구조를 분석하고
PX4/ArduPilot의 PPID 계단식 제어기와 EKF 상태 추정기를 MuJoCo 시뮬레이션으로 재현하여,
물리적 IMU 특성화에서 도출된 시뮬레이션 노이즈 조건 하에서 안정적인 자세 추적을 입증하였다.
논의 섹션에서는 IMU 노이즈와 지터가 내부 루프 각속도 추정과 PPID 추적 성능에 직접 영향을 미침을 보임으로써
두 연구를 연결하고 비행 컨트롤러 설계에서의 필터 선택과 이득 튜닝에 대한 실용적 지침을 제시하였다.

향후 연구에서는 하드웨어 측 정량적 평가를 위한 모션 캡처 시스템을 도입하고,
SK450 물리적 파라미터로부터 MuJoCo 모델을 직접 도출하며,
PX4 EKF2 아키텍처와 일치하는 완전한 23상태 EKF를 포함하도록 필터 비교를 확장할 예정이다.
